\section{The general calendar and the proper calendars}
\label{sec:particular}

\subsection{Two kinds of calendar}

\begin{normtext}{\loc{nualc} 48}
Ordinatio celebrationis anni liturgici calendario regitur, quod est generale vel particulare, prouti
statutum est in usum totius Ritus romani, vel in usum alicuius Ecclesiæ particularis vel familiæ
religiosæ.
\end{normtext}

\begin{normtext}{\loc{nualc} 49}
In calendario generali universus cyclus celebrationum inscribitur, tum mysterii salutis in Proprio de
tempore, tum eorum Sanctorum qui momentum universale præ se ferunt, et ideo obligatorie ab omnibus
celebrantur, tum aliorum qui universalitatem et continuitatem sanctitatis in populo Dei demonstrant.\\
Calendaria vero particularia celebrationes magis proprias continent congruenti ratione cum cyclo
generali organice compositas. Singulæ enim Ecclesiæ vel familiæ religiosæ Sanctos, qui ipsis sunt
peculiari ratione propria, peculiari honore prosequantur oportet.\\
Calendaria tamen particularia, a competenti auctoritate componenda, ab Apostolica Sede sunt probanda.
\end{normtext}

Three classes are distinguished inside the general calendar itself and they correspond exactly to
the three grades of the sanctoral: the mystery of salvation in the Proper of Time; those saints
whose significance is universal, celebrated obligatorily by all; and others who show the
universality and continuity of holiness in the people of God --- which is the theological description
of what an optional memorial is for. The particular calendar contains what is more properly its own,
organically composed with the general cycle, and it requires the approval of the Apostolic See.

\subsection{Composition: a particular calendar is the general calendar plus insertions}

\begin{normtext}{\loc{nualc} 52, opening}
Calendarium particulare conficitur inserendo calendario generali sollemnitates, festa et memorias
quæ sunt propria \dots
\end{normtext}

The verb is \latin{inserere}. A particular calendar is not an independent document beside the General
Roman Calendar and it is not a list of local additions; it is the general calendar with proper
celebrations inserted into it. The instruction \latin{Calendaria particularia} of 24~June 1970 says
the same thing in its own words at n.~13 --- \latin{Calendarium particulare conficitur, cum in
Calendarium generale celebrationes particulares inseruntur} --- and then enumerates the levels at
which such a calendar exists: national or regional, diocesan, religious.

The layering is cumulative and its direction is fixed. The 1970 instruction, n.~15\,b, composes a
diocesan calendar by inserting into the general calendar both the celebrations proper or granted to
the whole nation, region or wider territory \emph{and} those proper or granted to the whole diocese;
and n.~15\,c composes the calendar of each place, church, oratory and of religious congregations
without a calendar of their own from that diocesan calendar, adding their proper and indult
celebrations. So a parish church in the United States stands at the end of a chain of four:

\begin{ruletable}
\textbf{General Roman Calendar} & The universal cycle, in the typical edition of the Missal.\\
\textbf{National calendar} & The Proper Calendar for the Dioceses of the United States of America,
approved by the conference of bishops and confirmed by the Apostolic See, inserted into the
general calendar.\\
\textbf{Diocesan calendar} & The national calendar plus the diocese's own: its principal patron as a
feast (or, for pastoral reasons, a solemnity), the anniversary of the dedication of the cathedral as
a feast, its secondary patron as a memorial, and the saints and blessed connected with it by origin,
long residence, death, or immemorial and persisting cult (1970 instruction, n.~9).\\
\textbf{The church's own calendar} & The diocesan calendar plus the celebrations proper to that
church: the anniversary of its dedication as a solemnity if it is consecrated, its title as a
solemnity, and a memorial of the saint or blessed whose body is preserved there (1970 instruction,
n.~11; \loc{nualc}~52\,c).\\
\end{ruletable}

This is why the table of liturgical days assigns proper solemnities (place~4) and proper feasts
(place~8) their own graded sub-lists. The internal order within place~4 --- principal patron of the
place, dedication and its anniversary of the church itself, title of the church itself, then the
religious title, founder or principal patron --- is the resolution rule for collisions among the
proper layers, and it is the only place where the general law adjudicates between a diocesan and a
church claim.

\subsection{Restraints on a particular calendar}

The Norms and the 1970 instruction are noticeably more concerned with limiting particular calendars
than with enlarging them, and the restraints are worth listing because they explain why a diocesan
calendar is short.

\begin{ruletable}
\textbf{The Proper of Time is untouchable} & \loc{nualc}~50\,a: the Proper of Time is always to be
preserved entire and to enjoy its due pre-eminence over particular celebrations. The 1970 instruction
n.~2\,a adds the corollary: no perpetual particular celebration may be assigned to a Sunday.\\
\textbf{One celebration per saint} & \loc{nualc}~50\,b: individual saints are to enjoy a single
celebration in the liturgical year, with a second permitted as an optional memorial, where pastoral
reasons suggest it, for a translation or finding of patrons or founders.\\
\textbf{Indult celebrations are restrained} & \loc{nualc}~50\,c: celebrations granted by indult are
not to duplicate others already occurring in the cycle of the mystery of salvation, nor to be
multiplied beyond measure.\\
\textbf{Protected days stay clear} & \loc{nualc}~56\,f: so that the cycle of the liturgical year may
shine in its full light and the celebrations of saints not be perpetually impeded, the days on which
Lent and the Octave of Easter usually fall, and the days from 17 to 31 December, are to remain free
of particular celebrations, except optional memorials, the feasts of table n.~8\,a--d, and
solemnities which cannot be transferred to another time.\\
\textbf{Crowded calendars are pruned} & \loc{nualc}~53: where a diocese or religious family is
adorned with many saints and blessed, a common celebration of them all may be had; individual
celebration is reserved to those of particular significance for the whole diocese or family; the
rest are celebrated only in the places most closely connected with them or where their bodies are
kept.\\
\textbf{The default grade is low} & \loc{nualc}~54: proper celebrations are to be inscribed as
obligatory or optional memorials unless the table of liturgical days provides otherwise for some of
them or particular historical or pastoral reasons exist. The 1970 instruction, n.~24, gives the
reason: an optional memorial, precisely because it leaves the choice open, does not impede the
celebration of the saints but allows the day to be ordered to the needs and character of the
participants.\\
\end{ruletable}

Once approved, the calendar binds and freezes:

\begin{normtext}{\loc{nualc} 55}
Celebrationes calendario proprio inscriptæ ab omnibus qui ad illud calendarium tenentur servari
debent; et nonnisi approbante Apostolica Sede e calendario expungi vel gradu mutari possunt.
\end{normtext}

Neither removal nor a change of grade is within the local authority's own competence. This is the
structural reason why a national calendar changes slowly and why an annual calendar published by a
conference is a witness to the recurring calendar rather than an amendment of it.

\subsection{The proper day of a celebration}

\loc{nualc}~56 governs the date at which a proper celebration is inscribed, and its provisions are
about calendar construction, not about resolving a particular year:

\begin{itemize}
\item the Church's custom is to celebrate saints on their \latin{dies natalicius}, and this is to be
kept in proper celebrations too (56, opening);
\item a celebration also found in the general calendar is inscribed on the same day, with the grade
changed if necessary (56\,a);
\item a celebration not in the general calendar is assigned to the \latin{dies natalicius}; if that
day is unknown, to some other day proper to the saint --- ordination, finding, translation --- and
otherwise to a day free in the particular calendar (56\,b);
\item if the proper day is impeded by another obligatory celebration, even of lower grade, in the
general or particular calendar, the celebration is assigned to the nearest day similarly unimpeded
(56\,c);
\item but if the celebration cannot for pastoral reasons be moved, the impeding celebration is moved
instead (56\,d);
\item indult celebrations are inscribed on the day more suitable pastorally (56\,e).
\end{itemize}

Provision 56\,c is the one most often misapplied. It is a rule for assigning a permanent date when a
calendar is drawn up, and its trigger is a \emph{permanent} impediment on a fixed date, not the
occurrence of a movable celebration in one year. Nothing in it authorises moving a memorial in a
given year because Easter fell early.

\subsection{The United States layer}

The Proper Calendar for the Dioceses of the United States of America is the national layer for the
2011 English edition of the Missal. This reference does not reproduce it; the complete recurring
inventory of universal and national entries is a separate undertaking, and it is deliberately outside
this leaf's scope, which is the norms and the temporal cycle. What belongs here is the way the
national layer engages the rules stated above, and four kinds of engagement can be illustrated from
the Secretariat of Divine Worship's own published calendars.

\subsubsection{Assignment of solemnities to Sunday under \loc{nualc}~7}

The Body and Blood of Christ is assigned to the Sunday after Trinity throughout the United States;
the Secretariat's calendars print it as such, with the marker \textsc{usa}. The Epiphany is likewise
assigned to the Sunday occurring from 2 to 8 January. The Ascension is the interesting case, because
the assignment is not national but provincial: for 2026 the Secretariat records that Thursday 14~May
is the solemnity in the ecclesiastical provinces of Boston, Hartford, New York, Omaha and
Philadelphia, and that all other provinces of the United States have transferred it to the following
Sunday, 17~May, keeping the Feast of Saint Matthias on the Thursday. One country, one national
calendar, two different liturgical days on one civil date. No computation can produce this; it must
be read from the competent decision.

\subsubsection{Days of obligation, which are canonical and not calendrical}

The same calendars record that the obligation to attend Mass on 15~August 2026 is abrogated because
the date falls on a Saturday, under the complementary norm to canon 1246~\S2 confirmed for the
dioceses of the United States; and that by a decree of the Bishop of Honolulu dated 23~March 1992,
implementing a 1990 indult, the Nativity and the Immaculate Conception are the only two days of
obligation in the State of Hawaii. Neither statement changes a rank. The Assumption remains a
solemnity in table place~3 whether or not it is of precept, and the Hawaiian dioceses keep the same
solemnities as the rest of the country. The obligation and the calendar are separate systems that
touch only at \loc{nualc}~7, where the \emph{absence} of an obligation is the condition for
assigning a solemnity to Sunday.

\subsubsection{Civil observances admitted by inscription}

Independence Day and Thanksgiving Day appear in the Secretariat's calendar, and it explains exactly
why: they are inscribed in the Proper Calendar for the Dioceses of the United States of America,
whereas other civic holidays commonly kept in the country --- it names Martin Luther King Jr. Day,
Memorial Day and Labor Day --- are not, and the calendar therefore makes no suggestion for them
although the Missal supplies suitable Masses under \loc{igmr}~368--378. This is the clean
illustration of the difference between a celebration in a calendar and an occasion for which a Mass
may be chosen.

\subsubsection{What the annual calendar is, and is not}

The Secretariat's calendar is a resolution of one year, published so that authors of ordines and
other aids may work from a common text. It is the competent territorial witness to that resolution,
and this reference has used it as such. It is not the promulgating source of the recurring proper
calendar, and its own introduction points the reader elsewhere for the norms: to the Universal Norms,
particularly Chapter~II, for the composition of diocesan and provincial calendars, and to the General
Instruction and the Introduction to the Lectionary for the choice of texts. A recurring rule should
be taken from the books; a year's resolution should be checked against the calendar; and where the
two disagree, the disagreement is recorded and not resolved by preference.
